\section*{Hoofdstuk \ref{ch:b3:asymmetric-synthesis} --- Asymmetrische synthese}

\begin{solution}{exo:b3:asymmetric-synthesis:1}
ee 80\,\%: $\mathrm{er} = 1.8/0.2 = 9$ (90\,:\,10). er 99\,:\,1: $\mathrm{ee} = 98/100 = 98$\,\%. ee 50\,\%: 75\,\% en 25\,\%.
\end{solution}

\begin{solution}{exo:b3:asymmetric-synthesis:2}
$(1840 - 60)/(1840 + 60) = 0.937$: ee 93.7\,\%, een er van $1840/60 = 31$.
\end{solution}

\begin{solution}{exo:b3:asymmetric-synthesis:3}
Prioriteiten O $>$ ethyl $>$ H. Getekend met O boven, ethyl linksonder en H
rechtsonder is het voorvlak \textit{Si}, het achtervlak \textit{Re}. Een methylgroep die vanaf het \omterm{def:b3:asymmetric-synthesis:faces}{Re-vlak}
toegevoegd wordt, met de prioriteiten OH $>$ ethyl $>$ methyl $>$ H in het product, geeft (\textit{R})-butan-2-ol.
\end{solution}

\begin{solution}{exo:b3:asymmetric-synthesis:4}
Ethanol: \omterm{def:b3:asymmetric-synthesis:topicity}{enantiotoop} (verwisseld door het \omterm{def:b3:point-groups:mirror}{spiegelvlak} van het molecuul;
dezelfde verschuiving in een achiraal oplosmiddel). Butan-2-ol:
\omterm{def:b3:asymmetric-synthesis:topicity}{diastereotoop} (naast een stereocentrum; verschillende verschuivingen).
\end{solution}

\begin{solution}{exo:b3:asymmetric-synthesis:5}
er 199: $RT\ln 199 = \qty{13.1}{kJ/mol}$ bij \qty{298}{K}, \qty{8.6}{kJ/mol} bij \qty{195}{K}. 80\,\% ee bij \qty{25}{\celsius} is er 9,
$\Delta\Delta G^\ddagger = RT\ln 9 = \qty{5.45}{kJ/mol}$; bij \qty{-78}{\celsius} is $\mathrm{er} = \eu^{5450/(8.314 \times 195.15)} = 28.7$, ee 93\,\%.
\end{solution}

\begin{solution}{exo:b3:asymmetric-synthesis:6}
\omterm{def:b3:asymmetric-synthesis:ee}{Optische zuiverheid} $12.0/15.0 = 80$\,\%: 90\,\% (\textit{S}) en 10\,\% (\textit{R}). Een kleine draaiing,
onzuiverheden, concentratie- en oplosmiddeleffecten, en een
niet-lineair verband tussen draaiing en samenstelling maken ze
onbetrouwbaar; chromatografie scheidt en telt de enantiomeren
rechtstreeks.
\end{solution}

\begin{solution}{exo:b3:asymmetric-synthesis:7}
$s = \ln(0.40 \times 0.10)/\ln(0.40 \times 1.90) = -3.22/-0.274 = 12$.
\end{solution}

\begin{solution}{exo:b3:asymmetric-synthesis:8}
Met L-($+$)-diethyltartraat (2\textit{S},3\textit{S})-2,3-epoxyhexaan-1-ol; met D-($-$)-diethyltartraat de
enantiomeer, (2\textit{R},3\textit{R}).
\end{solution}

\begin{solution}{exo:b3:asymmetric-synthesis:9}
L = fenyl, M = methyl, S = H. Voor het aldehyde (\textit{S}) heeft de reactieve conformeer de fenylgroep
loodrecht op het carbonylvlak en het waterstofatoom naast de weg van het
nucleofiel, dat anti ten opzichte van fenyl aanvalt. Het uitwerken van de
configuraties geeft
(2\textit{S},3\textit{S})-3-fenylbutan-2-ol als belangrijkste diastereomeer.
\end{solution}

\begin{solution}{exo:b3:asymmetric-synthesis:10}
In het begin zijn de twee enantiomeren in gelijke hoeveelheden aanwezig,
zodat de producten ontstaan in de verhouding $k_{\mathrm{snel}} : k_{\mathrm{traag}} = s$: $\mathrm{ee}_p = (s - 1)/(s + 1)$, 0.905 voor $s = 20$. Naarmate de reactie vordert,
raakt de snelle enantiomeer uitgeput en daalt de ee van het product,
terwijl de ee van het substraat naar 1 stijgt: het substraat kan altijd
gezuiverd worden door verder te gaan; het product niet.
\end{solution}

\begin{solution}{exo:b3:asymmetric-synthesis:11}
Met snelle racemisatie: tot 100\,\% rendement, $\mathrm{er} = 99$, ee 98\,\%. Zonder racemisatie, bij 50\,\%
omzetting: een rendement van hoogstens 50\,\%, en een ee van het product
van 93\,\% (de ee van het substraat is dezelfde, 93\,\%).
\end{solution}

\begin{solution}{exo:b3:asymmetric-synthesis:12}
Als de twee complexen sneller in elkaar overgaan dan ze reageren, is de
productverhouding
$k_B[\mathrm B]/(k_A[\mathrm A]) = 600/10 = 60$: het minoritaire complex B geeft 98\,\% van het product (ee 97\,\%). De
selectiviteit wordt bepaald door het verschil tussen de twee
overgangstoestanden, niet door de populaties van de intermediairen.
\end{solution}

\begin{solution}{pb:b3:asymmetric-synthesis:1}
\textbf{1.} Het alkeenkoolstofatoom dat het stikstofatoom draagt, heeft drie
verschillende substituenten en twee verschillende vlakken; waterstof die
eraan toegevoegd wordt, maakt van het $\alpha$-koolstofatoom van het aminozuur een
stereocentrum.

\textbf{2.} Het racemaat: beide vlakken worden even snel aangevallen.

\textbf{3.} Zijn carbonylzuurstofatoom bindt samen met de binding C=C aan
het rodium: het substraat wordt als chelaat vastgehouden, in een vaste
geometrie die het ligand kan onderscheiden.

\textbf{4.} Het bouwt een chirale holte rond het metaal (symmetrie $C_2$): de twee manieren om
het substraat te binden, met het ene of het andere vlak, zijn
diastereomeer, met verschillende energieën en reactiviteiten.

\textbf{5.} Nee: het minoritaire complex reageert veel sneller met
waterstof (Curtin-Hammett).

\textbf{6.} Het metaal zit aan één vlak van het gebonden alkeen, en de
waterstofatomen worden vanaf het metaal overgedragen.

\textbf{7.} $\mathrm{er} = 1.95/0.05 = 39$.

\textbf{8.} $\Delta\Delta G^\ddagger = RT\ln 39 = 8.314 \times 298.15 \times 3.664 = \qty{9.1}{kJ/mol}$.

\textbf{9.} $\mathrm{er} = \eu^{9082/(8.314 \times 195.15)} = 270$: ee 99.3\,\%.

\textbf{10.} $\mathrm{er} = \eu^{11082/(8.314 \times 298.15)} = 87$: ee 97.7\,\%.

\textbf{11.} Minder dan een typische waterstofbrug: een klein verschil in
de sterische contacten of één zwakke interactie tussen substraat en
ligand beslist over de uitkomst, en daarom zijn zulke katalysatoren
moeilijk vanuit eerste beginselen te ontwerpen.

\textbf{12.} De snelheid daalt, waterstof lost anders op en de katalysator
gedraagt zich anders bij lage temperatuur, en een grote reactor koelen
kost energie.

\textbf{13.} 97.5 g van de enantiomeer (\textit{S}) en 2.5 g van de enantiomeer (\textit{R}).

\textbf{14.} $97.5 - 5.0 = \qty{92.5}{g}$ kristallen, praktisch enantiozuiver.

\textbf{15.} $92.5/100 = 92.5$\,\% van het ruwe product (95\,\% van de enantiomeer (\textit{S})).

\textbf{16.} De kristallen nemen de overmaat enantiomeer op, terwijl het
racemische deel met de minoritaire enantiomeer in oplossing blijft; uit
een racemaat valt geen overmaat te verzamelen, en de kristallisatie geeft
racemische kristallen.

\textbf{17.} Met chirale HPLC tegenover een racemische referentie.

\textbf{18.} 50\,\%: alleen de trage enantiomeer blijft over.

\textbf{19.} $s = \ln(0.45 \times 0.05)/\ln(0.45 \times 1.95) = -3.79/-0.131 = 29$.

\textbf{20.} 45\,\% van het racemaat (88\,\% van de gewenste enantiomeer die
in het begin aanwezig was, bij 95\,\% ee).

\textbf{21.} Met een snelle racemisatie van het substraat zou het rendement
100\,\% kunnen naderen.

\textbf{22.} Ze zet al de goedkope achirale precursor om in de juiste
enantiomeer, met een katalysator in zeer kleine hoeveelheden en zonder
enantiomeer weg te gooien.

\textbf{23.} Rodium is giftig en zijn resten in geneesmiddelen zijn beperkt
tot enkele ppm; ze worden gemeten met massaspectrometrie met inductief
gekoppeld plasma.

\textbf{24.} 95\,\% ee bij \qty{25}{\celsius} komt overeen met $\Delta\Delta G^\ddagger = RT\ln 39 \approx \qty{9.1}{kJ/mol}$.
\end{solution}
